\section{Introduction}\label{sec:introduction}

For nonnegative integer vectors $r=(r_1,\ldots,r_m)$ and
$c=(c_1,\ldots,c_n)$ with common total $N$, write
\[
 \Omega(r,c)=\left\{X\in\Z_{\ge0}^{m\times n}:
       \sum_jX_{ij}=r_i,\quad \sum_iX_{ij}=c_j\right\}.
\]
The uniform distribution on this finite set gives every table equal mass.
We study sampling when both dimensions vary and the margins are encoded
in binary. In particular, a polynomial bound in the numeric total $N$
would not, in general, give a polynomial bound in the input length.
There are no cell bounds or forbidden positions: every nonnegative integer
matrix with the prescribed margins is allowed. For probability laws on a
finite set, we use the convention
$\TV(\mu,\nu)=\tfrac12\sum_x|\mu(x)-\nu(x)|$.

\begin{theorem}\label{thm:main}
For arbitrary nonnegative integer margins of equal total, the following
algorithms exist in the model of unbiased random bits and bit operations.
\begin{enumerate}[label=\textup{(\roman*)}]
 \item Given an integer $k\ge1$, a bounded-time algorithm outputs a table
 in $\Omega(r,c)$ whose law is within $2^{-k}$ in total variation of the
 uniform law. Its running time is polynomial in
 $m,n,\log(N+1),k$.
 \item An algorithm outputs an exactly uniform table in $\Omega(r,c)$,
 terminates almost surely, and has expected running time polynomial in
 $m,n,\log(N+1)$.
\end{enumerate}
The polynomials are uniform over all margin vectors. No positivity,
sparsity, balance, or fixed-dimension hypothesis is imposed.
\end{theorem}

For a rational tolerance $0<\eps<1$ encoded by a binary numerator and
denominator, integer comparisons compute
$k=\lceil\log_2(1/\eps)\rceil$ in time polynomial in that encoding length.
Thereafter part~\textup{(i)} has accuracy dependence polynomial in
$\log(1/\eps)$.
The algorithms below are explicit, although their exponents are deliberately
large. Part~\textup{(ii)} concerns expected time: an exceptionally rare
branch performs exhaustive computation.

\paragraph{Counting and sampling fixed-margin arrays.}
Contingency tables are a basic family of integer arrays with linear
constraints. Their enumeration and exact and approximate generation are
treated, for example, by Diaconis and Gangolli~\cite{DiaconisGangolli1995}.
The algorithmic question involves more than producing one feasible table:
the law must give all feasible tables nearly or exactly the same mass,
even when their number and the coordinate ranges are exponentially large
in the margin bit length.

For two rows, Dyer and Greenhill~\cite{DyerGreenhill2000} give polynomial-time
approximate counting and sampling. Their heat-bath chain has mixing time
at most $n(n-1)\log(N/\eps)/2$ in the author version's Theorem~4.1.
Dependence on $\log N$ rather than $N$ makes this a polynomial bound in
the binary input length.
Kijima and Matsui~\cite[Theorems~3.1 and~5.1]{KM2003} give exact uniform
sampling for two rows by monotone coupling from the past, after sorting
the columns. Their expected $O(n^3\ln N)$ bound is stated for pseudocode
using uniform real draws. These results allow arbitrary two-row margin
sizes after deleting zero margins.

The fixed-row case extends beyond two rows. Cryan and
Dyer~\cite{CD2003} combine dynamic programming with volume estimation for
approximate counting and sampling. Dyer~\cite[Section~3]{Dyer2003} gives a
dynamic-programming rejection sampler that is exactly uniform, analyzed
in arithmetic operations. Cryan, Dyer, Goldberg,
Jerrum, and Martin~\cite{CDGJM2006} prove rapid mixing of the
$2\times2$ heat-bath chain for every fixed number of rows.
The polynomial bounds in these results depend on that fixed row count;
they do not give a bound polynomial jointly in both dimensions.

DeSalvo and Zhao~\cite[Theorem~2.3 and Remarks~2.4--2.5]{DZ2016} give a
general exact divide-and-conquer construction using $O(mn\log M)$ expected
random bits, where $M$ is the largest margin. In the cited arXiv version,
their polynomial total-runtime conclusion is conditional on efficient
evaluation of the parity-restricted counting quantities in Conjecture~2.6;
the bit-time bounds in Theorem~\ref{thm:main} are unconditional.

Another line of work uses the geometry of the transportation polytope
when margins are large. Dyer, Kannan, and Mount~\cite{DKM1997} develop
scaled coordinates, softened constraints, lattice walks, and rounding
comparisons for dense tables. Morris~\cite{Morris2002} improves the dense
sampling bounds through expanded polytopes and rounding. These methods
are close antecedents of the dense completion routine below. Our use of
that routine is conditional: a padding construction makes every
completion problem dense, while a separate chain handles all cells
incident to small margins.

For sparse margins, Arman, Gao, and Wormald~\cite[Theorem~1 and
Remark~2]{AGW2021} give exact uniform generation under $5\Delta^4<N$,
where $\Delta$ is the largest margin, in expected $O(N)$ time with
unit-cost arithmetic on integers of magnitude $O(N)$. The condition is
a genuine restriction, although within this family it also bounds $N$
polynomially in the dimensions. Their 2021 introduction surveys fixed-row,
dense, and sparse algorithms and reports that no polynomial-time
approximately uniform sampler for arbitrary margins was then known.
Theorem~\ref{thm:main} removes those margin and dimension restrictions
and includes all random bits and integer arithmetic in its cost.

\paragraph{Approximation and exact correction.}
Exactness and bounded running time are separate issues. Our approximate
algorithm has a deterministic polynomial work bound for each accuracy
request. The final exact correction uses the residual-mixture construction
of G\"obel, Liu, Manurangsi, and Pappik~\cite[Section~3.1, Algorithm~1 and
Theorem~6]{GLMP2024}. That construction combines a sufficiently accurate
sampler with a rare exhaustive draw from the remaining probability mass.
Its application here requires an additional property: the actual
implemented output law can be tabulated in time with an exponential
input-size factor but only polynomial dependence on the accuracy bits.
We prove that property and give the dyadic implementation explicitly.
Thus the exact algorithm has polynomial expected bit cost, not a
polynomial bound on every execution.

\paragraph{Proof strategy.}
Only cells incident to a margin below a polynomial threshold are explicit
coordinates of the auxiliary chain. The remaining cells form a complete
rectangle. We pad this rectangle so every conditional completion problem
has large margins. Each explicit cell has two bounded integer values,
one for its row and one for its column. Besides states where these values
agree, called transversals, we allow a positive unit discrepancy at one
cell and a negative unit discrepancy at another.

The target law on these small states assigns weight equal to the number
of padded completions. Sampling a state from that law and then a uniform
completion makes all state/completion pairs equally likely. Accept only
a transversal whose large completion entries are all at least the
padding. Subtracting it then gives a uniform original table conditional
on acceptance.
Section~\ref{sec:graph} constructs a repair map for the defect states and
uses it to bound their weight and the probability of successful unpadding.
The bounded sampler never evaluates completion counts: an ideal transition
uses a uniform completion draw to test a proposed move.

To prove that this chain mixes in polynomial time, we expose the small-cell
values of a transversal in a fixed order. After fixing a prefix, we compare
the averages of any real function on the state graph over the two fibers
that assign adjacent values to the next cell. The comparison comes from
quadratic forms indexed by pairs of unit removals.
Section~\ref{sec:signatures} proves that these forms have at most
one positive eigenvalue and that the property survives summing out a
balanced pair. A weighted telescoping calculation then bounds the energy
of a signed flow between neighboring conditional fibers.
Section~\ref{sec:transport} combines this bound with conditional-variance
decomposition and repair to obtain a Poincar\'e inequality for the
completion-count law. This supplies the small chain's mixing bound.

Section~\ref{sec:dense} implements the completion draws approximately on
a finite scaled grid, using rational arithmetic and dyadic random choices.
A continuous log-concave density serves only as an analytic tool for the
rejection and conductance bounds. Large coordinate ranges appear in
finite sums in the proof but are never enumerated in an ordinary run.
Section~\ref{sec:algorithms} couples each implemented completion draw to
its ideal counterpart and bounds the number of sampling and unpadding
attempts, giving the bounded almost-uniform algorithm. For exactness,
it also shows how to tabulate that algorithm's actual finite law,
including its stopping rules and fallbacks. This exhaustive computation
is performed only on the rare correction branch.
